The identity baseline returns the blended image unchanged. The threshold
baseline uses simple thresholding and connected-component logic to remove bright
regions. These baselines are intentionally simple references; thresholding can
perform worse than identity because removing bright pixels is not equivalent to
reconstructing hidden target structure.

Thayer-Direct follows the encoder-decoder U-Net design
\cite{ronneberger2015unet} and maps a blended image $X$ to a clean target image
$Y$:
\[
    \hat{Y} = f_{\theta}(X).
\]
The model is a compact encoder-decoder U-Net with skip connections.

Thayer-Residual predicts
\[
    R = X - Y,
\]
and reconstructs the target by subtraction:
\[
    \hat{Y} = X - g_{\theta}(X).
\]
This objective can preserve unchanged target light through the subtraction path.

Thayer-BR v0.1 is not a new architecture; it is a residual U-Net trained on a
more targeted blend distribution: 50\% normal/random blends, 30\%
high-overlap/core-obstruction blends, and 20\% brightness/size stress blends.
This tests whether targeted hard-case sampling improves robustness without
changing the basic model family.

Thayer-BR v0.2 Moderate keeps the residual U-Net formulation and balanced
training distribution target, but changes the loss. It uses a normalized
weighted residual mean squared error:
\[
    \mathcal{L} =
    \frac{\sum_i w_i \left(g_{\theta}(X)_i - (X_i - Y_i)\right)^2}
         {\sum_i w_i}.
\]
The affected mask is computed from the known synthetic pair,
\[
    \mathrm{mean}_c(|X - Y|) > 0.02,
\]
not from model predictions. The moderate setting adds weight to affected pixels
and additional weight to affected target-core pixels. Thayer-BR v0.2 Strong
uses larger affected/core weights and is treated as an ablation.
